\exnumber{00}
\extitle{Projet 1 - Get a room !}
\exfiles{ex00/archive zip, nomme Name\_FirstName\_DeliverableName\_Date\_VersionNumber}
\exauthorize{Ceux nommés dans l'exercice}
\makeheaderfiles

Création d'un service de réservation de salles sur desktop

Barbara travaille chez Coworking Square, un espace de coworking à Amsterdam.
Elle a besoin d'aide pour gérer les réservations. Parmi les 50 places de l'espace de travail, 35 peuvent être louées au mois ou à la semaine, et 15 sont des hot spots réservables pour une journée.

Pour l'instant, elle gère tout avec des calendriers papier mais ce n'est pas vraiment efficace.
« J'aimerais que mes clients puissent réserver des places à la semaine, au mois ou à la journée en ligne et tout payer d'avance pour que je n'aie pas à m'occuper de la facturation avant tout »
Elle vous appelle pour concevoir un site web qui pourrait résoudre son problème.

Vous devez la rencontrer après une semaine pour une première présentation de votre travail, après avoir réalisé un prototype mid-fidelity testé.
La présentation finale de votre travail aura lieu une semaine après la présentation intermédiaire, où vous montrerez un prototype haute fidélité testé.

\textbf{Durée estimée :} 2 semaines

\textbf{Taille de l'équipe :} 1

\textbf{Consignes :}

\begin{itemize}
\item Posez une stratégie UX et UI avant de foncer. Remettez le brief initial en question avec une recherche utilisateur.
\item Respectez les étapes : empathie, définir, énoncé du problème, idéation, prototype, test.
\item Vous devez interroger de vraies personnes pendant l'empathie, la recherche et les tests.
\item Dans les deux phases, testez vos prototypes avec au moins 5 personnes.
\item Faire une présentation comme si vous montriez votre travail à votre client. Assurez-vous qu'il y a au moins 5 personnes dans l'audience.
\item Déposez les documents sur Github, dans un zip nommé « ProjectTitle\_Name\_FirstName », et chaque livrable ainsi : « Name\_FirstName\_DeliverableName\_Date\_VersionNumber ». Exemple : « Doe\_John\_MidFiPrototype\_05242024\_V1 ».
\item Ajoutez un titre dans tous les documents écrits.
\item Soignez l'aspect global des livrables. Ils doivent être clairs, simples et lisibles. Vous pouvez regarder canva.com, sans surcharger le dessin de détails.
\item Préparez un formulaire de retour pour le public. Il contient au moins les questions listées plus bas.
\item Dans les deux présentations, vous parlez à Barbara et à son équipe.
\item Durée : 10 minutes de présentation et 5 minutes de questions.
\item Amusez-vous !
\end{itemize}

Outils :

\begin{itemize}
\item Google Forms \url{https://docs.google.com/forms/u/0/}
\item Airtable \url{https://www.airtable.com/}
\item Typeform \url{https://www.typeform.com/}
\end{itemize}

\textbf{Pistes :}

\begin{itemize}
\item Compétences de présentation efficaces \url{https://academic.oup.com/femsle/article/364/24/fnx235/4587905}
\end{itemize}

Livrables intermédiaires :

\textbf{Citation :}

« Your mission in any presentation is to inform, educate, and inspire. » Robert Ballard

\medskip

Phase 1 : Audit de la présentation mi-projet

Vérifiez :

\begin{itemize}
\item Au moins 5 personnes assistent à la présentation
\item La présentation porte sur tout le processus UX, y compris :
\begin{itemize}
\item Le brief initial
\item Un planning, avec les phases UX, les outils choisis, le calendrier
\item Les résultats de la recherche utilisateur, qualitative et quantitative (persona, chiffres, etc.)
\item Un énoncé de problème qui ajuste le brief initial d'après la recherche
\item Une introduction du concept sorti de l'idéation
\item Le prototype moyenne fidélité animé (une vidéo de démo)
\item Les retours principaux du test d'utilisabilité
\item Les prochaines étapes du projet
\end{itemize}
\item La présentation ne dépasse pas 10 minutes
\item Les formulaires de retour ont été distribués au début de la présentation
\item Vous vous êtes amusés !
\end{itemize}

\medskip

Phase 2 : Audit de la présentation finale

Vérifiez :

\begin{itemize}
\item Au moins 5 personnes assistent à la présentation
\item La présentation porte sur le processus UI, y compris :
\begin{itemize}
\item Un bref rappel de la présentation précédente
\item Un planning, avec les phases UI, les outils choisis, le calendrier
\item Une planche d'ambiance
\item Un styleguide (couleurs, marque, typographie, quelques éléments de la bibliothèque)
\item Le prototype haute fidélité animé (une vidéo de démo)
\item Les retours principaux des tests d'utilisabilité et de désirabilité
\item Les prochaines étapes du projet
\end{itemize}
\item La présentation ne dépasse pas 10 minutes
\item Les formulaires de retour ont été distribués au début de la présentation
\item Vous vous êtes amusés !
\end{itemize}

Questions à poser dans les formulaires de feedback

\begin{itemize}
\item Que retenez-vous de cette présentation ?
\item De 0 à 10, le discours était-il clair ?
\item Pourquoi ?
\item De 0 à 10, la présentation était-elle engageante ?
\item Pourquoi ?
\item De 0 à 10, la présentation était-elle cohérente ?
\item Pourquoi ?
\item De 0 à 10, les visuels portaient-ils ?
\item Pourquoi ?
\item De 0 à 10, le problème était-il clair ?
\item Pourquoi ?
\item De 0 à 10, le prototype répond-il au problème ?
\item Pourquoi ?
\item Quelle recommandation feriez-vous pour améliorer cette présentation ?
\item Pourquoi ?
\end{itemize}

